\section{Discussion}

\subsection{The frontier is the deliverable, not a defect}

H1 and H2 together describe a real precision/recall frontier, and the
practical upshot for an operator is that there is no policy choice that
dominates on every axis. Conservative (context-exposure) propagation is
the right default when the cost of a missed exposure exceeds the cost
of an unnecessary quarantine --- a compliance escalation, a legal hold,
an incident where under-flagging is the failure mode that ends
careers. Attribution thresholding is the right choice when quarantine
itself is expensive (every flagged object means a regenerated write, a
paused workflow, a human review) and some false-negative risk is
acceptable --- but H2's finding that an aggressive threshold (0.85)
specifically prunes the weakly-attributed \texttt{CO\_RETRIEVED} edges
that are sometimes the \emph{only} path to real contamination means
that choice should be made deliberately, not as a default tuning knob
turned up for cleanliness. \texttt{depth\_aware} (H4) is a third point
on this frontier, not a replacement for either: it recovers much of
conservative propagation's near-root recall while cutting long-tail
over-quarantine by roughly 30\%, at a missed-contamination cost
concentrated beyond depth 2. None of these is ``the'' answer; which one
an operator should run depends on what a false positive costs them
relative to a false negative, a judgment this paper does not make on
the operator's behalf.

Two robustness checks (\S\ref{sec:results}) address the obvious
``post-hoc operationalization'' criticism directly rather than leaving
it unaddressed: the attribution frontier's qualitative shape survives
under a second, architecturally different embedding model, and
\texttt{depth\_aware}'s cost/recall tradeoff moves smoothly and
monotonically as its conservative window is swept across $\{1, 2, 3\}$
hops rather than jumping around a single hand-picked point.

\subsection{No model is uniformly safer --- provenance metadata needs to know which model wrote what}

The cross-model asymmetry that recurs across H1, H3, and H4
(gpt-4o-mini shows the steepest inflation growth and the largest
\texttt{depth\_aware} missed-contamination cost; Llama-3.3-70B shows by
far the most laundering under \texttt{paraphrase}) is not just a
comparison result --- it has a direct operational consequence. A
deployment that mixes models (a cheap model for routine derivation, a
stronger model for high-stakes writes, a fallback during an outage)
should not assume one propagation policy tuned against a single
model's profile performs equally well against another's --- our
results motivate recording which model produced a given derived
memory and evaluating containment policy as a function of that model
identity, rather than assuming a uniform policy is safe by default. A
\texttt{depth\_aware} window tuned against gpt-4o's
missed-contamination profile is not obviously the same policy for a
memory Llama-3.3-70B wrote, though we did not test whether a uniform
policy is in fact acceptable in any particular deployment --- only that
per-model behavior is heterogeneous enough that the question is worth
asking. We did not evaluate policy parameters conditioned on model
identity; this paper
establishes that the asymmetry is real and large enough to matter, not
what the model-conditioned policy should be.

\subsection{The compositional boundary is real, not a labeler artifact}

The REFERENCES-class weakness flagged in Limitations (precision 0.75,
recall 0.46, concentrated in \texttt{multi\_hop\_setup} scenarios)
could have been explained two ways: genuine rubric ambiguity at the
CARRIES/REFERENCES boundary, or a labeler-specific failure. The
human-human $\kappa$ baseline ($\kappa = 0.9277$, 115/120 raw
agreement) supports the former: all five disagreements between two
independent human annotators fell on that same boundary, though two
annotators over 120 examples is evidence for, not a settled ontology
of, it. A second signal agrees: the manual, full-population read of
all 46 \texttt{child\_2=CARRIES} cases in the real-document slice
(\S\ref{sec:results}) found 40/46 explicitly contained the injected
claim's marker or wording---the labeler's CARRIES calls are not
pattern-matching noise on held-out data either. The operational
implication is that any deployment using this framework (or any
automated labeler at all) should treat multi-hop/compositional
exposure flags as lower-confidence by construction and route them to
human review, rather than trusting the same confidence threshold that
works for single-premise claims.

\subsection{Where this sits relative to prevention and reactive investigation}

Related Work positions this paper against MemLineage (preventive
gating) and MemAudit (reactive, triggered by an observed harmful event)
as answering a different question, not a better one. Read together,
the three form a plausible pipeline rather than competing solutions:
preventive lineage enforcement may prevent a compromised label from
ever reaching a sensitive action; reactive investigation triggers once
harm is observed and identifies which memory caused \emph{that
specific} failure; post-incident reconstruction --- this paper's object
--- is what an operator needs the moment either signal fires but
before its full downstream consequence is known, a question none of
the three papers this compares against evaluates over a branching
graph with a measured precision/recall/inflation frontier.

\subsection{When this framework is, and is not, the right tool}

The metrics and policies here assume an agent architecture where
derivation provenance (which memories were retrieved into which write)
can be logged at generation time --- true of retrieval-augmented agents
built on an explicit vector store or structured memory log, the
architecture this paper's threat model assumes. An agent whose memory
is an opaque fine-tuned weight update, or whose retrieval/write path is
not instrumented, cannot use this framework at all; blast-radius
reconstruction has to be designed in before an incident, not
retrofitted after one. We are not arguing every agent memory system
should adopt these specific metrics, only that systems capable of
logging provenance in the first place have a measurable
precision/recall/cost frontier to reason about, and, to our knowledge,
this paper is the first measurement of what that frontier looks like
under branching derivation.

\subsection{What this changes about where to look next}

The human-human baseline redirects the next step for the
REFERENCES-class weakness away from ``improve the labeler'' and toward
a targeted human study of that boundary, and motivates extending the
real-document slice (\S\ref{sec:real-document-slice}) to mixed-model
traces.
